\section{補助量子ビットと逆演算}
\label{sec:ancilla-uncomputation}

量子回路では，途中の計算結果を保持するために追加の量子ビットを用いる．
これを補助量子ビット，複数まとめたものを補助レジスタという．
補助レジスタに途中の情報が残ると，目的のレジスタと結び付いたままになる．
その結果，後の計算が正しく行えないことがある．
そこで，必要な操作を終えた後に計算を逆向きに行い，
補助レジスタを初期状態へ戻す．この操作を逆演算またはアンコンピュテーションという．

系レジスタの計算基底状態$\ket{x}$に対して
\begin{equation}
 U\ket{x}\ket0=\ket{x}\ket{f(x)}
 \label{eq:ancilla-computation}
\end{equation}
となるユニタリ演算$U$を考える．
次に，$f(x)$の値に応じて変わる制御演算$V$を別のレジスタに加える．
その後に$U^\dagger$を作用させると
\begin{align}
 \sum_x\alpha_x\ket{x}\ket0\ket\eta
 &\xmapsto{U}
 \sum_x\alpha_x\ket{x}\ket{f(x)}\ket\eta\\
 &\xmapsto{V}
 \sum_x\alpha_x\ket{x}\ket{f(x)}V_{f(x)}\ket\eta\\
 &\xmapsto{U^\dagger}
 \sum_x\alpha_x\ket{x}\ket0V_{f(x)}\ket\eta
 \label{eq:uncomputation-superposition}
\end{align}
となる．最後の等号は，$U^\dagger U=I$と線形性による．
標的の状態$V_{f(x)}\ket\eta$には$f(x)$に応じた変化が残る．
一方，計算用レジスタは$\ket0$に戻る．

ただし，$U^\dagger$を作用させる前に$\ket{x}$または$\ket{f(x)}$を測定したり，
式\eqref{eq:ancilla-computation}の対応を壊す操作を加えたりすると，
上の逆演算は成立しない．HHLでは，固有値に応じた制御回転を終えた後，
逆量子位相推定によって位相レジスタを初期状態へ戻す．
